\documentclass[11pt]{article}
\usepackage[utf8]{inputenc}
\usepackage[T1]{fontenc}
\usepackage[margin=2.5cm]{geometry}
\usepackage{amsmath,amssymb}
\usepackage{graphicx}
\usepackage{booktabs}
\usepackage{xcolor}
\usepackage[hidelinks]{hyperref}
\usepackage{authblk}
\usepackage{setspace}
\onehalfspacing

% ---------------------------------------------------------------------------
% Every result comes from generated/numbers.tex (scripts/analyze.py).
% \res{<model>}{<metric>} prints a value; missing values print a red ??.
% \tbd{...} marks text to be written once the final run is in.
% ---------------------------------------------------------------------------
\newcommand{\res}[2]{\ifcsname res@#1@#2\endcsname\csname res@#1@#2\endcsname\else\textcolor{red}{\textbf{??}}\fi}
\providecommand{\bestmodel}{\textcolor{red}{BEST}}
\providecommand{\worstmodel}{\textcolor{red}{WORST}}
\InputIfFileExists{generated/numbers.tex}{}{}
\newcommand{\tbd}[1]{\textcolor{red}{[#1]}}
\newcommand{\gentable}[1]{\IfFileExists{generated/tables/#1.tex}{\input{generated/tables/#1.tex}}{\begin{table}[ht]\centering\fbox{\textcolor{red}{Table \texttt{#1} is generated by scripts/analyze.py}}\end{table}}}
\newcommand{\genfigure}[3]{\begin{figure}[ht]\centering\IfFileExists{generated/figures/#1.pdf}{\includegraphics[width=#2\linewidth]{generated/figures/#1.pdf}}{\fbox{\textcolor{red}{Figure \texttt{#1} is generated by scripts/analyze.py}}}\caption{#3}\label{fig:#1}\end{figure}}

\title{Counting Photographs, Not Images: Duplicate-Aware Evaluation of Deep Learning for Skin Disease Classification}

% Continues and corrects the July 2026 preprint. Its co-authors are not
% continuing; a new collaborator can be added if they contribute substantially
% to this study and approve the final text.
\author{Farzana Rahman Supti}
\affil{Department of Electrical and Computer Engineering, North South University, Dhaka, Bangladesh\\ \texttt{farzanasupti@gmail.com}}
\date{}

\begin{document}
\maketitle

\begin{abstract}
Image-based classifiers for skin disease often report accuracies that do not survive independent evaluation. A common cause is duplicate or near-duplicate images that fall on both sides of a train--test split. We benchmarked a compact convolutional network and seven frozen ImageNet backbones on a public five-class dataset: acne, hyperpigmentation, nail psoriasis, Stevens--Johnson syndrome/toxic epidermal necrolysis and vitiligo (\res{data}{n_images} images). The images are \res{data}{original_images} photographs, each stored with 27 augmented copies, so a split by image places copies of most test photographs in training. The dataset has no patient identifiers, so images were grouped by source photograph, using the augmentation tool's filenames and perceptual hashes, before splitting, and only original photographs were scored. Models were compared by five-fold grouped cross-validation, and the best model and a three-model ensemble were chosen on validation folds only. On a held-out test set of \res{data}{n_test_images} photographs, the best single model (\bestmodel) reached \res{\bestmodel}{test_accuracy}\% accuracy (95\% CI \res{\bestmodel}{test_accuracy_ci}) and \res{\bestmodel}{test_f1_macro}\% macro-F1. The ensemble reached \res{Ensemble}{test_accuracy}\% accuracy (\res{Ensemble}{test_accuracy_ci}). No architecture was significantly better than another across folds, and no pair of models, the ensemble included, differed significantly on the test set. With Monte Carlo dropout, referring the 20\% most uncertain cases raised the ensemble's accuracy on the retained cases from \res{Ensemble}{acc_refer_0}\% to \res{Ensemble}{acc_refer_20}\%. MC dropout did not improve on the ensemble's own probabilities: its expected calibration error was \res{Ensemble}{ece_mc}\%, against \res{Ensemble}{ece_det}\% without it, and plain softmax entropy flagged errors as well as MC-dropout entropy. Code, splits and all predictions are public.
\end{abstract}

\noindent\textbf{Keywords:} skin disease classification; deep learning; data leakage; ensemble; Monte Carlo dropout; uncertainty; calibration

% ===========================================================================
\section*{Introduction}
Skin conditions are among the most common causes of non-fatal disease burden worldwide\cite{karimkhani2017gbd}. Many patients have no timely access to a dermatologist. Deep convolutional networks have matched specialists on curated dermoscopic tasks\cite{esteva2017,haenssle2018} and on differential diagnosis from clinical photographs\cite{liu2020}. Performance often falls, however, when models meet images from new sources or patient populations\cite{daneshjou2022}. This gap between reported and real-world accuracy has several causes. One of the most avoidable is information leaking from the test set into training\cite{kapoor2023,varoquaux2022,roberts2021}.

In image classification, leakage most often arises when copies or near-copies of an image appear in both the training and the test partitions. Splitting at the patient level prevents this when patient identifiers exist. In brain MRI, slice-level rather than subject-level splitting inflated reported accuracy substantially\cite{yagis2021}. Public dermatology collections assembled from several sources also contain duplicates. An analysis of the ISIC archives found duplicated images across the training and test sets of past challenges\cite{cassidy2022}. Datasets gathered partly from online sources, as many disease-specific collections are, carry the same risk and often have no patient metadata with which to remove it.

A second gap is uncertainty. A model that is accurate on average is clinically useful only if it can recognise the cases it is likely to get wrong and refer them to a clinician\cite{leibig2017}. Monte Carlo (MC) dropout approximates Bayesian inference by keeping dropout active at test time\cite{gal2016dropout,kendall2017}. Ensembles of independently trained networks are a strong, simple alternative\cite{lakshminarayanan2017,ovadia2019}, and they are also widely used to improve skin-lesion classification accuracy\cite{harangi2018}. Neither is useful unless its uncertainty is checked: how well calibrated the probabilities are\cite{guo2017calibration}, and whether uncertain predictions really are the wrong ones.

We study a public five-class dataset of acne, hyperpigmentation, nail psoriasis, Stevens--Johnson syndrome/toxic epidermal necrolysis (SJS-TEN) and vitiligo\cite{khushbu2024dataset}. These conditions are under-represented in the melanoma-focused benchmarks that dominate the field. Our contributions are:
\begin{enumerate}
  \item a duplicate-aware evaluation protocol for datasets without patient identifiers, which groups images by source photograph (from filenames and perceptual hashes) before grouped, stratified splitting, scores original photographs only, and reports how many duplicates it finds;
  \item a comparison of a compact custom CNN with seven frozen ImageNet backbones under five-fold grouped cross-validation, with each backbone given the input preprocessing its weights expect, and with significance testing across folds and on the test set;
  \item a cross-architecture ensemble whose members are selected on validation folds only, evaluated on a held-out test set with group-bootstrap confidence intervals;
  \item an evaluation of MC-dropout uncertainty for calibration, error detection and uncertainty-based referral.
\end{enumerate}
All code, the exact splits and every model prediction are released so that each number in this paper can be regenerated. This study supersedes an earlier preprint on the same dataset, whose results were affected by the duplicate problem described here (see \emph{Relation to an earlier preprint}).

% ===========================================================================
\section*{Results}

\subsection*{Duplicates in the dataset}
The dataset contains \res{data}{n_images} images in \res{data}{n_classes} classes, but only \res{data}{original_images} of them are original photographs. The other \res{data}{augmented_copies} are augmented copies made before release: each photograph appears with 27 copies, cropped, padded, rendered as superpixels, shifted in hue or saturation, flipped and resized, and named after the photograph it came from. Grouping by source filename linked \res{data}{links_from_filenames} pairs of images, and perceptual hashing linked another \res{data}{links_from_hashes} (Methods), \res{data}{links_from_flips_rotations} of them only after mirroring or rotating one image of the pair. \res{data}{border_cropped_images} images had a uniform frame that was removed before hashing. Altogether \res{data}{images_in_multi_image_groups} images fell into \res{data}{n_multi_image_groups} multi-image groups (largest group: \res{data}{max_group_size} images), leaving \res{data}{n_groups} independent groups. \res{data}{groups_with_conflicting_labels} group contained images with different class labels. The duplicate problem is therefore structural rather than marginal: 96.4\% of the images are copies of another image in the dataset, and the effective sample is \res{data}{original_images} photographs, not \res{data}{n_images} images. Under an image-level split that held out one image in seven, the chance that none of a test photograph's 27 copies falls into training is $(1/7)^{27}<10^{-22}$, so in practice every test image would have copies in training. The group with conflicting labels is also the largest one. Hash links joined the copies of one vitiligo photograph to those of four acne photographs; these are different photographs of similar appearance, not one photograph filed under two labels. The group was kept whole, in development fold~1. Every image kept its original label, and the group does not touch the test set. The held-out test set contains \res{data}{n_test_images} original photographs (\res{data}{n_test_images_with_copies} images with their copies), and the five development folds together contain \res{data}{n_dev_images} (\res{data}{n_dev_images_with_copies}). No group spans two partitions.

\subsection*{Model comparison under cross-validation}
Table~\ref{tab:cv} and Figure~\ref{fig:cv_folds} show validation performance across the five folds. Mean macro-F1 ranked \bestmodel{} first and \worstmodel{} last (\res{\worstmodel}{cv_f1_macro}\% $\pm$ \res{\worstmodel}{cv_f1_macro_sd}\%). A Friedman test across folds did not reject equal performance of all models ($\chi^2_{\res{summary}{friedman_df}} = \res{summary}{friedman_chi2}$, $p = \res{summary}{friedman_p}$). The Nemenyi critical difference was \res{summary}{nemenyi_cd} rank units. That is larger than the gap between any two mean ranks, which ran from 2.4 (\bestmodel) to 5.6 (\worstmodel), so no pair of models is separated. % mean ranks from summary.json
Mean validation macro-F1 spanned \res{\worstmodel}{cv_f1_macro}--\res{\bestmodel}{cv_f1_macro}\%, and its fold-to-fold standard deviations (2.7--6.3 points) were of the same order as the differences between models. The compact CNN, trained from scratch, reached \res{CustomCNN}{cv_f1_macro}\% $\pm$ \res{CustomCNN}{cv_f1_macro_sd}\% macro-F1. That is inside the range of the frozen pretrained backbones, though below the best three.

\gentable{cv}
\genfigure{cv_folds}{0.7}{Validation macro-F1 of each model on each of the five development folds. Horizontal bars: means.}

\subsection*{Held-out test performance}
For each architecture, the five fold models were averaged and evaluated once on the held-out test set (Table~\ref{tab:test}). \bestmodel{} reached \res{\bestmodel}{test_accuracy}\% accuracy (95\% CI \res{\bestmodel}{test_accuracy_ci}), \res{\bestmodel}{test_f1_macro}\% macro-F1 and an MCC of \res{\bestmodel}{test_mcc}. The ensemble of the \res{summary}{n_members} best architectures on cross-validation (\res{summary}{members}) reached \res{Ensemble}{test_accuracy}\% accuracy (\res{Ensemble}{test_accuracy_ci}) and \res{Ensemble}{test_f1_macro}\% macro-F1. Against \bestmodel{} alone, an exact McNemar test gave a Holm-adjusted $p = \res{summary}{mcnemar_ens_vs_best_p}$. Over all model pairs, \res{summary}{n_sig_pairs} of \res{summary}{n_pairs} comparisons remained significant after Holm correction (Supplementary Table~\ref{tab:mcnemar}). The ensemble classified 40 of the 45 test photographs correctly and \bestmodel{} 39.

Test results agreed with the cross-validation estimates: \bestmodel{} reached \res{\bestmodel}{test_accuracy}\% on test, against \res{\bestmodel}{cv_accuracy}\% $\pm$ \res{\bestmodel}{cv_accuracy_sd}\% across folds. Single-architecture test accuracy ranged from \res{InceptionV3}{test_accuracy}\% (InceptionV3) to \res{EfficientNetB0}{test_accuracy}\% (EfficientNetB0). With \res{data}{n_test_images} test photographs, one photograph moves accuracy by 2.2 points, and all models' confidence intervals overlap. The order on the test set also differed from the order in cross-validation: the compact CNN, second to last across folds, matched \bestmodel{} on test (\res{CustomCNN}{test_accuracy}\%). Rankings are unstable at this sample size.

The ensemble misclassified 5 of the 45 test photographs (Table~\ref{tab:perclass-Ensemble}, Figure~\ref{fig:confusion}). All nail psoriasis photographs and all three hyperpigmentation photographs were classified correctly. Two of 15 SJS-TEN photographs were called acne. Three of 10 vitiligo photographs were also misclassified, two as SJS-TEN and one as acne. Acne therefore had full recall but the lowest precision (62.5\%). \bestmodel{} alone also missed two SJS-TEN and three vitiligo photographs, and one of the three hyperpigmentation photographs. With three test photographs, the perfect hyperpigmentation result says little. % recheck error counts after regeneration

\gentable{test}
\gentable{per_class_Ensemble}
\genfigure{confusion}{0.95}{Held-out test confusion matrices for the best single architecture (left) and the ensemble (right). Cells give counts; shading gives the row-normalised rate.}

\subsection*{Calibration and uncertainty}
Table~\ref{tab:uncertainty} compares deterministic predictions with predictions averaged over \res{data}{mc_samples} MC-dropout samples per fold model. For the ensemble, expected calibration error was \res{Ensemble}{ece_det}\% without and \res{Ensemble}{ece_mc}\% with MC dropout (Figure~\ref{fig:reliability}). Used to separate wrong from correct predictions, MC-dropout predictive entropy gave an AUROC of \res{Ensemble}{err_auroc_mc_entropy}. Softmax entropy of the deterministic ensemble gave \res{Ensemble}{err_auroc_softmax_entropy}, and mutual information gave \res{Ensemble}{err_auroc_mutual_info}. Referring the most uncertain cases raised accuracy on the retained cases from \res{Ensemble}{acc_refer_0}\% (no referral) to \res{Ensemble}{acc_refer_10}\% at 10\% referral and \res{Ensemble}{acc_refer_20}\% at 20\% (Figure~\ref{fig:referral}).

MC dropout added nothing measurable to the ensemble. Its calibration error was slightly higher with MC dropout than without, and its negative log-likelihood was almost unchanged (Table~\ref{tab:uncertainty}). For error detection, the deterministic ensemble's softmax entropy did as well as MC-dropout entropy. For single models the picture was mixed. MC dropout lowered calibration error for five of eight architectures (for \bestmodel, from \res{\bestmodel}{ece_det}\% to \res{\bestmodel}{ece_mc}\%) but raised it for VGG16, InceptionV3 and Xception. With \res{data}{n_test_images} test photographs, each calibration bin holds only a few photographs, so these differences are small relative to their own uncertainty. Referral helped modestly. Referring 20\% of cases (9 photographs) removed two of the ensemble's five errors. Referring 30\% (14 photographs) removed all five, which is why accuracy on the retained cases reached \res{Ensemble}{acc_refer_30}\%. Both misclassified SJS-TEN photographs fell between the 20\% and 30\% thresholds (uncertainty ranks 10 and 12 of 45). % recheck counts and ranks after regeneration

\gentable{uncertainty}
\genfigure{reliability}{0.5}{Reliability diagram on the held-out test set. A perfectly calibrated model lies on the diagonal.}
\genfigure{referral}{0.55}{Accuracy on retained cases as the most uncertain cases (by MC-dropout predictive entropy) are referred. Dotted: referring cases at random.}

\subsection*{Model size and cost}
Table~\ref{tab:efficiency} lists parameter counts, training time and inference time. The custom CNN has \res{CustomCNN}{params_trainable} trainable parameters. The frozen backbones train only their classifier heads but still carry their full ImageNet weights at inference. In absolute terms the cost differences were small: every model classified an image in 1--2.3\,ms in batches on the GPU. \tbd{Xception's training and inference times include fold~4, which was run on a CPU after the Kaggle GPU quota ran out; replace once it is rerun on the T4.} The compact CNN, with \res{CustomCNN}{params_total_m}\,M parameters in total, was the smallest model and the fastest at inference (\res{CustomCNN}{ms_batched}\,ms per image). It matched \bestmodel{} (\res{\bestmodel}{params_total_m}\,M) on the test set, but it ranked lower across folds and took the longest to train, because every layer is trained. Among the backbones, EfficientNetB0 (\res{EfficientNetB0}{params_total_m}\,M) gave the best balance: second across folds, the highest single-model test accuracy, and one of the fastest to train.

\gentable{efficiency}

% ===========================================================================
\section*{Discussion}
The main finding concerns the data rather than the models. The dataset's \res{data}{n_images} images reduce to \res{data}{original_images} independent photographs, because each photograph was released with 27 augmented copies. A split by image would place copies of essentially every test photograph in training, so a model could score well by recognising photographs it had already seen rather than by recognising disease. Grouping by source photograph removes that route, and it shrinks the test set to \res{data}{n_test_images} photographs. The price is wide confidence intervals: the ensemble's 95\% interval for accuracy spans \res{Ensemble}{test_accuracy_ci}\%. We think these intervals, not the point estimates, are the honest summary of what this dataset can show. Datasets released in pre-augmented form should state so, and should keep the link between each copy and its source, so that users can split by source.

Under this protocol the architectures were statistically indistinguishable. No pair was separated across folds, and no pair differed on the test set after correction. The compact CNN, trained from scratch at $128\times128$ pixels, fell inside the range of the frozen ImageNet backbones. It was below the best three across folds and level with the best on test. Pretrained features therefore gave no clear advantage at this resolution and data size. Fine-tuning the backbones, or training at higher resolution, might change this, and we did not test either.

The ensemble was comparable to the best single model, not better: it was right on one more test photograph (McNemar $p=\res{summary}{mcnemar_ens_vs_best_p}$). Its practical gains were elsewhere. It had the lowest negative log-likelihood and Brier score of all models, and its uncertainty separated errors from correct predictions well (AUROC \res{Ensemble}{err_auroc_softmax_entropy} from softmax entropy alone). MC dropout added nothing on top of the ensemble, for calibration, error detection or referral. This null result is consistent with reports that ensembles capture most of the predictive uncertainty that dropout sampling provides\cite{lakshminarayanan2017,ovadia2019}. It is also expected here, because dropout acts only on the classifier heads of frozen backbones. For a single model, MC dropout lowered calibration error in five of eight architectures. An ensemble of a few architectures is the simpler route to the same end.

The errors that matter most clinically are the two missed SJS-TEN photographs, both classified as acne by the ensemble. SJS-TEN is a dermatological emergency, and labelling it as acne would delay treatment. Uncertainty-based referral would have caught both photographs only at 30\% referral, not at 20\%. A system meant for triage would therefore need either a higher referral rate or a decision threshold that favours SJS-TEN. Vitiligo was the other weak class, with three of ten photographs misclassified. The hyperpigmentation results rest on three test photographs and should not be interpreted.

\paragraph{Limitations.} The dataset has no patient identifiers. Grouping by filename and perceptual hash removes exact and near-exact copies, but it cannot ensure that different photographs of the same patient stay in one partition. Our estimates may therefore still be optimistic for truly unseen patients. The dataset holds only \res{data}{original_images} independent photographs, so the test set has \res{data}{n_test_images} and the smallest class, hyperpigmentation, has only a handful; confidence intervals are correspondingly wide, and per-class results for the small classes are anecdotal. Perceptual hashing also joined some different photographs taken in similar settings (bare backs against a dark background, close-ups of nails); this keeps them in one partition and cannot leak, but it slightly reduces the number of independent groups. The perceptual hashes do not match copies that were cropped heavily or rotated by angles other than multiples of $90^\circ$. Everything was evaluated on one dataset, collected from hospitals and online sources with unknown acquisition protocols, demographics and skin-tone distribution. There is no external test set, and performance on other populations\cite{daneshjou2022} is unknown. Images were downsampled to $128\times128$ pixels and the pretrained backbones were not fine-tuned, both of which may understate what these architectures can achieve. MC dropout acts only on the classifier heads, so it captures uncertainty in those layers and not in the frozen feature extractors. The system is a research benchmark and has not been evaluated prospectively or against clinicians.

\paragraph{Conclusion.} The \res{data}{n_images} images in this public dataset come from only \res{data}{original_images} photographs, and evaluation must be grouped by photograph to be valid. Evaluated that way, a compact CNN, frozen ImageNet backbones and a three-model ensemble reached test accuracies of \res{InceptionV3}{test_accuracy}--\res{Ensemble}{test_accuracy}\% with overlapping confidence intervals, and none was significantly better than another. MC dropout did not improve on the ensemble's own uncertainty. A reliable comparison will need more independent photographs, ideally with patient identifiers, and an external test set.

% ===========================================================================
\section*{Methods}

\subsection*{Dataset}
We used version 1 of the Skin Disease Classification Dataset\cite{khushbu2024dataset}. It has five classes: acne (1{,}148 images), hyperpigmentation (700), nail psoriasis (2{,}520), SJS-TEN (3{,}164) and vitiligo (2{,}016). The images were collected from hospitals and online sources, and the dataset provides no patient identifiers or other metadata. The release is already augmented: it holds \res{data}{original_images} original photographs, and each appears with 27 copies whose filenames record the operation and the photograph they came from (for example, \texttt{Crop-0350\_03003227} is a crop of photograph \texttt{3227}). We used the archive from Mendeley Data, in which each class is a nested zip file. For Kaggle, these were unpacked into one folder per class without changing any image, and all \res{data}{n_images} images were loaded. Images were read in RGB colour order and resized to $128\times128$ pixels with area interpolation. Raw pixel values in $[0,255]$ were passed to the models, each of which applies its own input normalisation (below). The dataset is public and de-identified, and no new data were collected, so ethics approval was not required.

\subsection*{Duplicate-aware grouping}
Each image was assigned to a group with a union--find procedure over two kinds of link. (i) \emph{Filename links}: images in the same class whose filenames reduce to the same source stem after removing the extension, export-tool suffixes of the form \texttt{\_jpg.rf.<hash>}, and the augmentation prefix (\texttt{Crop-}, \texttt{padd-}, \texttt{superpixels-}, \texttt{hue-}, \texttt{saturation-}, \texttt{flip-} or \texttt{Resize-} with its numeric parameters). Because the parameters run straight into the source name, a prefixed name was matched to the longest original filename in its class that it ends with. These links are essential: crops and superpixel renderings of a photograph do not resemble it closely enough for a perceptual hash, and an earlier run of this study that grouped by hashes alone left every one of the 341 photographs with copies on both sides of the train--test split. (ii) \emph{Hash links}: images that look the same after allowing for re-encoding, resizing, recolouring, padding, mirroring and rotation. Each image was converted to greyscale and any uniform frame around it was cropped, so that a copy shrunk onto a black canvas is compared by its content. Two 64-bit perceptual hashes\cite{zauner2010} were then computed for each of the eight flips and $90^\circ$ rotations of the cropped image: a difference hash and a discrete-cosine-transform hash. Two images were linked when, for the stored form of one and some flip or rotation of the other, both hashes differed in at most 8 of 64 bits. Requiring both hashes to agree is essential. In a pilot on this dataset, the difference hash alone linked unrelated black-padded and smooth images, and union--find chained those links into one group of 561 images spanning all five classes. Hash links are made across classes, so a copy filed under two labels is also caught. Linked images form one group, and a group is never split between partitions. We chose the threshold to err towards merging: a spurious link only keeps two unrelated images in the same partition, whereas a missed link lets a copy of a training image into evaluation.

\subsection*{Splits}
A stratified group split (scikit-learn \texttt{StratifiedGroupKFold}, seed 42) set aside one seventh of the groups (\res{data}{n_test_images} of the \res{data}{original_images} photographs, about 13\% of the images) as a held-out test set. The remaining development set was divided into five grouped, stratified folds. The test set was not used for any training, model selection, early stopping or hyperparameter choice. Models were trained, and stopped early, on all images in their partitions, augmented copies included, but every validation and test metric counts original photographs only: the copies of one photograph are not independent cases, and scoring them would weight each photograph 28-fold and narrow confidence intervals without adding evidence. The split of every image is released with the code.

\subsection*{Models}
The custom CNN has four blocks of two $3\times3$ convolutions (16, 32, 64 and 192 filters), each followed by batch normalisation\cite{ioffe2015batchnorm} and a LeakyReLU activation (slope 0.1), with $2\times2$ max pooling after each block. Without batch normalisation, one of five folds in a pilot run never left the trivial solution. A head follows: global average pooling, a 1{,}536-unit dense layer with LeakyReLU, dropout (rate 0.4) and a five-way softmax. Its inputs are scaled to $[0,1]$.

The seven transfer models use ImageNet-pretrained\cite{deng2009imagenet} VGG16\cite{simonyan2015vgg}, ResNet50\cite{he2016resnet}, MobileNetV2\cite{sandler2018mobilenetv2}, EfficientNetB0\cite{tan2019efficientnet}, DenseNet121\cite{huang2017densenet}, InceptionV3\cite{szegedy2016inception} and Xception\cite{chollet2017xception} backbones without their top layers. The backbones are frozen and run in inference mode, so their batch-normalisation statistics are fixed\cite{tajbakhsh2016}. Each backbone receives the preprocessing its weights were trained with: RGB-to-BGR conversion and ImageNet mean subtraction for VGG16 and ResNet50; scaling to $[-1,1]$ for MobileNetV2, InceptionV3 and Xception; ImageNet mean and standard-deviation normalisation for DenseNet121; and raw $[0,255]$ input for EfficientNetB0, which rescales internally. Each backbone is followed by global average pooling and a trainable head of one or two ReLU dense layers (512, or 512 and 256 units) with dropout (0.3--0.5). Batch normalisation follows the first dense layer for ResNet50, EfficientNetB0 and Xception. The exact heads are in the released code. All dense layers carry an L2 weight penalty of $10^{-4}$.

\subsection*{Training}
Models were trained with Adam\cite{kingma2015adam} (learning rate $10^{-3}$, batch size 32) to minimise categorical cross-entropy. Class weights were inversely proportional to class frequency in the training folds. Augmentation was applied on the fly to training images only: random horizontal and vertical flips, rotation up to $\pm20^\circ$, translation up to $\pm20\%$ and zoom up to $\pm20\%$. Training ran for at most 150 epochs (custom CNN) or 50 epochs (transfer models). The learning rate was halved after 10 epochs without improvement in validation loss (minimum $10^{-7}$), and training stopped after 20 such epochs. The weights from the epoch with the lowest validation loss were kept. Each fold used its own fixed random seed, shared by all architectures so that they see the same augmentation stream. Experiments ran in Keras~3 with the TensorFlow backend on a Kaggle NVIDIA Tesla T4 GPU (16\,GB) with mixed precision.

\subsection*{Ensembling and model selection}
For each architecture, the softmax outputs of its five fold models were averaged to give one test prediction. Architectures were ranked by mean validation macro-F1 across folds. The best-ranked architecture is reported as the best single model, and the top three form the ensemble, again by averaging probabilities. The test set played no part in either choice.

\subsection*{Monte Carlo dropout and uncertainty}
At inference, dropout layers were kept active while batch normalisation and augmentation stayed in inference mode. Each fold model was sampled $T=\res{data}{mc_samples}$ times\cite{gal2016dropout}. Pooling the samples over fold models (and over ensemble members) gave a predictive distribution $\bar p$. From it we computed predictive entropy $H[\bar p]$ and mutual information $H[\bar p] - \frac{1}{S}\sum_s H[p_s]$, which is the epistemic part of the uncertainty\cite{kendall2017}. Calibration was measured by expected calibration error (15 equal-width confidence bins)\cite{guo2017calibration}, negative log-likelihood and the multi-class Brier score. Error detection was measured by the AUROC of each uncertainty score for separating misclassified from correctly classified test images. For referral, the test images were sorted by predictive entropy and accuracy was recomputed on the retained images as the most uncertain 5--30\% were removed\cite{leibig2017}.

\subsection*{Metrics and statistics}
We report accuracy, balanced accuracy, macro- and weighted-F1, Matthews correlation coefficient\cite{chicco2020mcc}, Cohen's $\kappa$ and macro one-vs-rest AUC. Test-set confidence intervals are 95\% percentile intervals from 2{,}000 bootstrap resamples of \emph{groups}, not images\cite{efron1993}. Cross-validation comparisons used the Friedman test on per-fold validation macro-F1, with mean ranks and the Nemenyi critical difference ($\alpha=0.05$)\cite{demsar2006}. With five folds these tests have limited power. Pairwise comparisons on the test set used exact McNemar tests on per-image correctness\cite{dietterich1998}, Holm-corrected over all pairs\cite{holm1979}. Reporting follows the TRIPOD+AI checklist where it applies\cite{collins2024tripod}.

% ===========================================================================
\section*{Data availability}
The dataset is publicly available from Mendeley Data (\href{https://doi.org/10.17632/3hckgznc67.1}{doi:10.17632/3hckgznc67.1}). The image-level split, duplicate groups and perceptual hashes used in this study are included in the code repository.

\section*{Code availability}
All code, the Kaggle notebook used to produce the results, every model's validation and test predictions, and the scripts that generate each table and figure are available at \url{https://github.com/farzanasupti/skin-disease-classification} \tbd{and archived at Zenodo, DOI}.

\section*{Relation to an earlier preprint}
This study continues and corrects an earlier preprint on the same dataset (ResearchGate, July 2026, doi:10.13140/RG.2.2.27428.03208), on which the present author was first author. \textbf{The results of that preprint are superseded and should not be cited as evidence of model performance.} Three problems affect them. First, the data were split by image. Because the dataset consists of \res{data}{original_images} photographs, each released with 27 augmented copies (Results), an image-level split places copies of nearly every test photograph in the training data, so the accuracies it reports are inflated. Second, its pretrained models received incorrectly preprocessed inputs. Third, its reported results could not be reproduced from its code. The present study redesigned and reran every experiment and rewrote the text; none of the preprint's results, figures or text are reused. The co-authors of the preprint did not take part in the present study.

\section*{Acknowledgements}
The author thanks the co-authors of the earlier preprint for their contributions to that work.

\section*{Author contributions}
F.R.S. conceived and designed the study, wrote the software, ran the experiments, analysed the results and wrote the manuscript.

\section*{Use of AI tools}
An AI assistant (Claude, Anthropic) was used to help write and test the code, to audit the dataset for duplicates, and to draft and edit the manuscript text. The author directed this work, checked all of the assistant's output against the code and the results, and takes full responsibility for the content. \tbd{Check this statement against the target journal's policy on AI-assisted writing.}

\section*{Competing interests}
The authors declare no competing interests.

\bibliographystyle{unsrt}
\bibliography{references}

\clearpage
\appendix
\section*{Supplementary information}
\gentable{mcnemar}
\gentable{per_class_\bestmodel}

\end{document}
